\subsection{Dark Energy as Emergent Projective Pressure}
  \label{subsec:absence-of-dark-energy-signatures}

  Because cosmic acceleration emerges as a geometric consequence of the
  global growth of projected capacity, no independent dark energy component is introduced.
  Vacuum energy is not absent from the framework: it arises as the residual
  projective capacity of configurations not yet stably locked by the projection ---
  the portion of the non-injective fibre that remains unresolved at each cascade
  step~(Section~\ref{subsec:vacuum-energy-relaxation-capacity}).
  What the framework does not introduce is a vacuum energy \emph{tuned independently}
  to explain cosmic acceleration, nor an additional propagating degree of freedom
  responsible for it.

  However, the absence of an independent dark energy \emph{component} does not
  imply a constant equation of state.
  The projective cascade exhibits two structurally distinct growth regimes
  (Section~\ref{sec:relax-cascade}): a polynomial regime governed by individual BFS
  shell growth ($|S_n| \sim n^3$, homogeneous dimension $D_{\mathrm{hom}} = 4$), and
  an exponential regime governed by the branching of distinct admissible trajectories.
  The effective equation-of-state parameter~$w(z)$ evolves across the transition
  between these two regimes: near $w < -1$ when exponential branching dominates at
  early times, relaxing toward $w > -1$ as the admissible neighbourhood
  $\mathcal{A}(U_n)$ becomes progressively constrained by accumulated projective
  history.
  This structural prediction is compatible with recent DESI~DR2 evidence
  favouring dynamical dark energy over a cosmological
  constant~\cite{DESI2025}.

  The key observational discriminators of the projective mechanism from
  standard dark energy models are therefore:
  \begin{itemize}
    \item no independent clustering component or fifth force associated with a
    dark energy field;
    \item a specific evolution of $w(z)$ tied to the polynomial-to-exponential
    transition of the admissible projective volume, rather than to a scalar potential;
    \item the combined absence of an inflationary tensor imprint at large angular
    scales and suppressed low-$\ell$ CMB power, as signatures of the same projective
    mechanism governing both early and late acceleration.
  \end{itemize}
